% Options for packages loaded elsewhere
\PassOptionsToPackage{unicode}{hyperref}
\PassOptionsToPackage{hyphens}{url}
\PassOptionsToPackage{dvipsnames,svgnames,x11names}{xcolor}
%
\documentclass[
]{article}
\usepackage{amsmath,amssymb}
\usepackage{iftex}
\ifPDFTeX
  \usepackage[T1]{fontenc}
  \usepackage[utf8]{inputenc}
  \usepackage{textcomp} % provide euro and other symbols
\else % if luatex or xetex
  \usepackage{unicode-math} % this also loads fontspec
  \defaultfontfeatures{Scale=MatchLowercase}
  \defaultfontfeatures[\rmfamily]{Ligatures=TeX,Scale=1}
\fi
\usepackage{lmodern}
\ifPDFTeX\else
  % xetex/luatex font selection
  \setmonofont[Scale=0.85]{DejaVu Sans Mono}
\fi
% Use upquote if available, for straight quotes in verbatim environments
\IfFileExists{upquote.sty}{\usepackage{upquote}}{}
\IfFileExists{microtype.sty}{% use microtype if available
  \usepackage[]{microtype}
  \UseMicrotypeSet[protrusion]{basicmath} % disable protrusion for tt fonts
}{}
\makeatletter
\@ifundefined{KOMAClassName}{% if non-KOMA class
  \IfFileExists{parskip.sty}{%
    \usepackage{parskip}
  }{% else
    \setlength{\parindent}{0pt}
    \setlength{\parskip}{6pt plus 2pt minus 1pt}}
}{% if KOMA class
  \KOMAoptions{parskip=half}}
\makeatother
\usepackage{xcolor}
\usepackage[a4paper,margin=2.2cm]{geometry}
\usepackage{color}
\usepackage{fancyvrb}
\newcommand{\VerbBar}{|}
\newcommand{\VERB}{\Verb[commandchars=\\\{\}]}
\DefineVerbatimEnvironment{Highlighting}{Verbatim}{commandchars=\\\{\}}
% Add ',fontsize=\small' for more characters per line
\newenvironment{Shaded}{}{}
\newcommand{\AlertTok}[1]{\textcolor[rgb]{1.00,0.00,0.00}{\textbf{#1}}}
\newcommand{\AnnotationTok}[1]{\textcolor[rgb]{0.38,0.63,0.69}{\textbf{\textit{#1}}}}
\newcommand{\AttributeTok}[1]{\textcolor[rgb]{0.49,0.56,0.16}{#1}}
\newcommand{\BaseNTok}[1]{\textcolor[rgb]{0.25,0.63,0.44}{#1}}
\newcommand{\BuiltInTok}[1]{\textcolor[rgb]{0.00,0.50,0.00}{#1}}
\newcommand{\CharTok}[1]{\textcolor[rgb]{0.25,0.44,0.63}{#1}}
\newcommand{\CommentTok}[1]{\textcolor[rgb]{0.38,0.63,0.69}{\textit{#1}}}
\newcommand{\CommentVarTok}[1]{\textcolor[rgb]{0.38,0.63,0.69}{\textbf{\textit{#1}}}}
\newcommand{\ConstantTok}[1]{\textcolor[rgb]{0.53,0.00,0.00}{#1}}
\newcommand{\ControlFlowTok}[1]{\textcolor[rgb]{0.00,0.44,0.13}{\textbf{#1}}}
\newcommand{\DataTypeTok}[1]{\textcolor[rgb]{0.56,0.13,0.00}{#1}}
\newcommand{\DecValTok}[1]{\textcolor[rgb]{0.25,0.63,0.44}{#1}}
\newcommand{\DocumentationTok}[1]{\textcolor[rgb]{0.73,0.13,0.13}{\textit{#1}}}
\newcommand{\ErrorTok}[1]{\textcolor[rgb]{1.00,0.00,0.00}{\textbf{#1}}}
\newcommand{\ExtensionTok}[1]{#1}
\newcommand{\FloatTok}[1]{\textcolor[rgb]{0.25,0.63,0.44}{#1}}
\newcommand{\FunctionTok}[1]{\textcolor[rgb]{0.02,0.16,0.49}{#1}}
\newcommand{\ImportTok}[1]{\textcolor[rgb]{0.00,0.50,0.00}{\textbf{#1}}}
\newcommand{\InformationTok}[1]{\textcolor[rgb]{0.38,0.63,0.69}{\textbf{\textit{#1}}}}
\newcommand{\KeywordTok}[1]{\textcolor[rgb]{0.00,0.44,0.13}{\textbf{#1}}}
\newcommand{\NormalTok}[1]{#1}
\newcommand{\OperatorTok}[1]{\textcolor[rgb]{0.40,0.40,0.40}{#1}}
\newcommand{\OtherTok}[1]{\textcolor[rgb]{0.00,0.44,0.13}{#1}}
\newcommand{\PreprocessorTok}[1]{\textcolor[rgb]{0.74,0.48,0.00}{#1}}
\newcommand{\RegionMarkerTok}[1]{#1}
\newcommand{\SpecialCharTok}[1]{\textcolor[rgb]{0.25,0.44,0.63}{#1}}
\newcommand{\SpecialStringTok}[1]{\textcolor[rgb]{0.73,0.40,0.53}{#1}}
\newcommand{\StringTok}[1]{\textcolor[rgb]{0.25,0.44,0.63}{#1}}
\newcommand{\VariableTok}[1]{\textcolor[rgb]{0.10,0.09,0.49}{#1}}
\newcommand{\VerbatimStringTok}[1]{\textcolor[rgb]{0.25,0.44,0.63}{#1}}
\newcommand{\WarningTok}[1]{\textcolor[rgb]{0.38,0.63,0.69}{\textbf{\textit{#1}}}}
\usepackage{longtable,booktabs,array}
\usepackage{calc} % for calculating minipage widths
% Correct order of tables after \paragraph or \subparagraph
\usepackage{etoolbox}
\makeatletter
\patchcmd\longtable{\par}{\if@noskipsec\mbox{}\fi\par}{}{}
\makeatother
% Allow footnotes in longtable head/foot
\IfFileExists{footnotehyper.sty}{\usepackage{footnotehyper}}{\usepackage{footnote}}
\makesavenoteenv{longtable}
\setlength{\emergencystretch}{3em} % prevent overfull lines
\providecommand{\tightlist}{%
  \setlength{\itemsep}{0pt}\setlength{\parskip}{0pt}}
\setcounter{secnumdepth}{5}
\ifLuaTeX
\usepackage[bidi=basic]{babel}
\else
\usepackage[bidi=default]{babel}
\fi
\babelprovide[main,import]{italian}
% get rid of language-specific shorthands (see #6817):
\let\LanguageShortHands\languageshorthands
\def\languageshorthands#1{}
\ifLuaTeX
  \usepackage{selnolig}  % disable illegal ligatures
\fi
\usepackage{bookmark}
\IfFileExists{xurl.sty}{\usepackage{xurl}}{} % add URL line breaks if available
\urlstyle{same}
\hypersetup{
  pdftitle={VERA X16 per Atari --- Eseguire i test nell'emulatore},
  pdflang={it},
  colorlinks=true,
  linkcolor={blue},
  filecolor={Maroon},
  citecolor={Blue},
  urlcolor={Blue},
  pdfcreator={LaTeX via pandoc}}

\title{VERA X16 per Atari --- Eseguire i test nell'emulatore}
\usepackage{etoolbox}
\makeatletter
\providecommand{\subtitle}[1]{% add subtitle to \maketitle
  \apptocmd{\@title}{\par {\large #1 \par}}{}{}
}
\makeatother
\subtitle{Fork di atari800 con la scheda VeraX16 PBI}
\author{Gianluca Renzi}
\date{2026-10-06}

\begin{document}
\maketitle

{
\hypersetup{linkcolor=}
\setcounter{tocdepth}{3}
\tableofcontents
}
\section{Panoramica}\label{panoramica}

I programmi di test in \texttt{vera-tests/} girano su un fork
dell'emulatore \textbf{atari800} (versione 5.2.0) che emula la scheda
VeraX16 sul bus PBI dell'Atari: registri VERA a \texttt{\$D100-\$D11F},
ROM handler da 2 KB a \texttt{\$D800-\$DFFF} (selezionata dal latch PBI
\texttt{\$D1FF}).

L'emulatore si trova accanto a questo repository:

\begin{verbatim}
~/Progetti/ATARI/
├── atari800/                 emulatore (src/atari800 è l'eseguibile)
│   └── vera_pbi_rom -> ../VERA_ATARI_PBI
└── VERA_ATARI_PBI/           questo repository (ROM, driver, test)
\end{verbatim}

Tutti i comandi seguenti si eseguono \textbf{dalla radice di questo
repository} (\texttt{VERA\_ATARI\_PBI}), quindi l'emulatore è
\texttt{../atari800/src/atari800}. Per brevità:

\begin{Shaded}
\begin{Highlighting}[]
\VariableTok{EMU}\OperatorTok{=}\NormalTok{../atari800/src/atari800}
\end{Highlighting}
\end{Shaded}

\section{Prerequisiti}\label{prerequisiti}

\subsection{Compilare l'emulatore}\label{compilare-lemulatore}

\begin{Shaded}
\begin{Highlighting}[]
\BuiltInTok{cd}\NormalTok{ ../atari800}
\ExtensionTok{./autogen.sh}
\ExtensionTok{./configure} \AttributeTok{{-}{-}enable{-}pbi{-}verax16}
\FunctionTok{make}
\BuiltInTok{cd} \AttributeTok{{-}}
\end{Highlighting}
\end{Shaded}

L'emulazione della VERA è in \texttt{src/pbi\_verax16.c} (registri, FX,
audio, SPI) e \texttt{src/vera\_video.c} (rendering). Senza
\texttt{-\/-enable-pbi-verax16} le opzioni \texttt{-verax16} non
esistono.

\subsection{ROM del sistema operativo
Atari}\label{rom-del-sistema-operativo-atari}

Se atari800 è compilato con libcurl (\texttt{configure} attiva da solo
la funzione di download quando la trova) e non è configurato nessun file
ROM del sistema operativo, all'avvio \textbf{scarica le ROM da
Internet}: preleva
\texttt{http://www.emulators.com/freefile/pcxf380.zip} (PC Xformer 3.8),
estrae i file \texttt{.rom} in \texttt{../atari800/src/rom/}, li usa (OS
XL e BASIC originali Atari) e ne scrive i percorsi in
\texttt{../atari800/src/.atari800.cfg}. Da quel momento non li scarica
più.

atari800 legge \texttt{.atari800.cfg} dalla cartella dell'eseguibile
prima di \texttt{\textasciitilde{}/.atari800.cfg}: una volta creato
\texttt{../atari800/src/.atari800.cfg}, lanciando
\texttt{../atari800/src/atari800} il file
\texttt{\textasciitilde{}/.atari800.cfg} non viene più letto.

Per usare AltirraOS e Altirra BASIC integrati, senza download, compilare
con:

\begin{Shaded}
\begin{Highlighting}[]
\ExtensionTok{./configure} \AttributeTok{{-}{-}enable{-}pbi{-}verax16} \AttributeTok{{-}{-}disable{-}download}
\end{Highlighting}
\end{Shaded}

e lasciare vuoti i percorsi \texttt{ROM\_OS\_*} nel file di
configurazione.

\subsection{Compilare ROM, driver e dischi di
test}\label{compilare-rom-driver-e-dischi-di-test}

Strumenti richiesti: \texttt{cc65} (\texttt{ca65}, \texttt{ld65},
\texttt{cl65}), \texttt{dir2atr}, \texttt{python3}.

\begin{Shaded}
\begin{Highlighting}[]
\FunctionTok{make}            \CommentTok{\# ROM, driver VERA*.SYS, programmi di test, immagini disco}
\end{Highlighting}
\end{Shaded}

Prodotti principali:

\begin{longtable}[]{@{}
  >{\raggedright\arraybackslash}p{(\columnwidth - 2\tabcolsep) * \real{0.5000}}
  >{\raggedright\arraybackslash}p{(\columnwidth - 2\tabcolsep) * \real{0.5000}}@{}}
\toprule\noalign{}
\begin{minipage}[b]{\linewidth}\raggedright
File
\end{minipage} & \begin{minipage}[b]{\linewidth}\raggedright
Contenuto
\end{minipage} \\
\midrule\noalign{}
\endhead
\bottomrule\noalign{}
\endlastfoot
\texttt{vera\_pbi\_handler.rom} & ROM handler PBI, 2 KB (avvio in
80×60) \\
\texttt{VERA4030.SYS}, \texttt{VERA8030.SYS}, \texttt{VERA8060.SYS} &
driver in RAM per 40×30, 80×30, 80×60 \\
\texttt{disk1-runcpm.atr} & DOS 2.0S, \texttt{RUNCPM.COM},
\texttt{VERA8030.SYS} \\
\texttt{disk2-veratests-40x30.atr} & \texttt{TEST4}, \texttt{TESTGS4},
\texttt{TESTMAZ4}, \texttt{TESTMTX4}, \texttt{TESTRMT},
\texttt{VERA4030.SYS} \\
\texttt{disk2-veratests-80x30.atr} & \texttt{TEST8}, \texttt{TESTGS8},
\texttt{TESTMAZ8}, \texttt{TESTMTX8}, \texttt{TESTRMT},
\texttt{VERA8030.SYS} \\
\texttt{disk2-veratests-80x60.atr} & \texttt{TEST6}, \texttt{TESTGS6},
\texttt{TESTMAZ6}, \texttt{TESTMTX6}, \texttt{TESTRMT},
\texttt{VERA8060.SYS} \\
\texttt{disk3-standalone.atr} & \texttt{TESTFX}, \texttt{TESTIRQ},
\texttt{TESTPLR} + \texttt{DEMO.VTM} (test senza driver) \\
\texttt{disk4-rmtio.atr} & \texttt{TESTRIO} con autorun MyPicoDos e i
file di prova \\
\end{longtable}

I programmi
\texttt{TESTn}/\texttt{TESTGSn}/\texttt{TESTMAZn}/\texttt{TESTMTXn} dei
dischi \texttt{disk2} contengono già il driver della loro risoluzione
(\texttt{n} = 4: 40×30, 8: 80×30, 6: 80×60).

\section{Il comando di base}\label{il-comando-di-base}

\begin{Shaded}
\begin{Highlighting}[]
\VariableTok{$EMU} \AttributeTok{{-}xl} \AttributeTok{{-}pal} \AttributeTok{{-}nopatch} \AttributeTok{{-}verax16} \AttributeTok{{-}verax16{-}rom}\NormalTok{ vera\_pbi\_handler.rom }\DataTypeTok{\textbackslash{}}
\NormalTok{     disk3{-}standalone.atr}
\end{Highlighting}
\end{Shaded}

\begin{itemize}
\tightlist
\item
  \texttt{-verax16} inserisce la scheda; \texttt{-verax16-rom} indica la
  ROM handler.
\item
  \texttt{-xl} (800XL) o \texttt{-xe} (130XE): il driver richiede una
  macchina XL/XE.
\item
  \texttt{-pal} o \texttt{-ntsc}: i test devono passare in entrambi.
\item
  \texttt{-nopatch}: nessuna patch SIO ad alta velocità, l'I/O su disco
  va alla velocità seriale reale (il dispositivo H: resta attivo). Per
  renderlo permanente impostare \texttt{ENABLE\_SIO\_PATCH=0} in
  \texttt{\textasciitilde{}/.atari800.cfg}.
\item
  L'ultimo argomento è il disco in \texttt{D1:}; altri \texttt{.atr}
  vanno in \texttt{D2:}, \texttt{D3:}\ldots{}
\end{itemize}

L'emulatore mostra due uscite: lo schermo ANTIC/GTIA normale e lo
schermo della VERA (VGA 640×480).

\subsection{Avviare un programma dal DOS
2.0S}\label{avviare-un-programma-dal-dos-2.0s}

I dischi avviano il DOS 2.0S. Dal menu del DUP:

\begin{enumerate}
\def\labelenumi{\arabic{enumi}.}
\tightlist
\item
  premere \texttt{L} (BINARY LOAD);
\item
  scrivere il nome del file, per esempio \texttt{TESTFX.COM}, e premere
  Return.
\end{enumerate}

\subsection{Avviare un programma senza
disco}\label{avviare-un-programma-senza-disco}

\texttt{-run} carica un eseguibile direttamente dal file system
dell'host:

\begin{Shaded}
\begin{Highlighting}[]
\VariableTok{$EMU} \AttributeTok{{-}xl} \AttributeTok{{-}pal} \AttributeTok{{-}nopatch} \AttributeTok{{-}verax16} \AttributeTok{{-}verax16{-}rom}\NormalTok{ vera\_pbi\_handler.rom }\DataTypeTok{\textbackslash{}}
     \AttributeTok{{-}run}\NormalTok{ TESTIRQ.COM}
\end{Highlighting}
\end{Shaded}

\section{Opzioni}\label{opzioni}

\subsection{Opzioni della scheda
VeraX16}\label{opzioni-della-scheda-verax16}

\begin{longtable}[]{@{}
  >{\raggedright\arraybackslash}p{(\columnwidth - 2\tabcolsep) * \real{0.5000}}
  >{\raggedright\arraybackslash}p{(\columnwidth - 2\tabcolsep) * \real{0.5000}}@{}}
\toprule\noalign{}
\begin{minipage}[b]{\linewidth}\raggedright
Opzione
\end{minipage} & \begin{minipage}[b]{\linewidth}\raggedright
Significato
\end{minipage} \\
\midrule\noalign{}
\endhead
\bottomrule\noalign{}
\endlastfoot
\texttt{-verax16} & attiva la scheda (alias \texttt{-\/-use-verax16}) \\
\texttt{-verax16-rom\ F} & ROM handler, 2 KB a \texttt{\$D800-\$DFFF} \\
\texttt{-verax16-pbi-id\ N} & bit del dispositivo PBI 0-7 (default 7,
maschera \texttt{\$80}) \\
\texttt{-verax16-config-ms\ N} & tempo di configurazione dell'FPGA dopo
l'accensione, bus non pilotato nel frattempo (default 100); \texttt{0} =
la scheda tiene il RESET dell'Atari fino a CONFIG\_DONE \\
\texttt{-verax16-debuglevel\ N} & livello di log (default 0) \\
\texttt{-verax16-psg-volume\ N} & livello del PSG VERA in percento,
0-400 (default 100 = una voce forte come un canale POKEY a volume 15) \\
\texttt{-verax16-sdcard\ F} & immagine SD grezza (per esempio da
\texttt{dd}) esposta tramite la SPI della VERA \\
\end{longtable}

\subsection{Opzioni utili di atari800}\label{opzioni-utili-di-atari800}

\begin{longtable}[]{@{}
  >{\raggedright\arraybackslash}p{(\columnwidth - 2\tabcolsep) * \real{0.5000}}
  >{\raggedright\arraybackslash}p{(\columnwidth - 2\tabcolsep) * \real{0.5000}}@{}}
\toprule\noalign{}
\begin{minipage}[b]{\linewidth}\raggedright
Opzione
\end{minipage} & \begin{minipage}[b]{\linewidth}\raggedright
Uso
\end{minipage} \\
\midrule\noalign{}
\endhead
\bottomrule\noalign{}
\endlastfoot
\texttt{-xl}, \texttt{-xe} & macchina 800XL / 130XE \\
\texttt{-pal}, \texttt{-ntsc} & standard video \\
\texttt{-basic}, \texttt{-nobasic} & ROM BASIC attiva/disattiva \\
\texttt{-run\ F} & esegue un file COM/EXE/XEX/BAS \\
\texttt{-nopatch} & nessuna patch SIO ad alta velocità: temporizzazione
seriale reale, H: funziona \\
\texttt{-nopatchall} & nessuna patch all'OS: temporizzazione seriale SIO
reale (H: non funziona) \\
\texttt{-volume\ N} & volume di uscita 0-100 \\
\texttt{-stereo} & due POKEY \\
\texttt{-turbo} & massima velocità \\
\texttt{-netsio\ {[}porta{]}} & NetSIO per FujiNet-PC (default UDP
9997) \\
\texttt{-config\ F} & file di configurazione alternativo \\
\end{longtable}

\section{Eseguire i singoli test}\label{eseguire-i-singoli-test}

\subsection{TESTFX --- coprocessore FX}\label{testfx-coprocessore-fx}

\begin{Shaded}
\begin{Highlighting}[]
\VariableTok{$EMU} \AttributeTok{{-}xl} \AttributeTok{{-}pal} \AttributeTok{{-}nopatch} \AttributeTok{{-}verax16} \AttributeTok{{-}verax16{-}rom}\NormalTok{ vera\_pbi\_handler.rom }\DataTypeTok{\textbackslash{}}
\NormalTok{     disk3{-}standalone.atr}
\end{Highlighting}
\end{Shaded}

Dal DUP: \texttt{L}, \texttt{TESTFX.COM}. Verifica ogni registro FX e
misura il throughput di copia/riempimento della VRAM. Risultato atteso:
\textbf{PASS 36, FAIL 0}.

\subsection{TESTIRQ --- interrupt della
VERA}\label{testirq-interrupt-della-vera}

\begin{Shaded}
\begin{Highlighting}[]
\VariableTok{$EMU} \AttributeTok{{-}xl} \AttributeTok{{-}pal}  \AttributeTok{{-}nopatch} \AttributeTok{{-}verax16} \AttributeTok{{-}verax16{-}rom}\NormalTok{ vera\_pbi\_handler.rom }\DataTypeTok{\textbackslash{}}
     \AttributeTok{{-}run}\NormalTok{ TESTIRQ.COM}
\VariableTok{$EMU} \AttributeTok{{-}xl} \AttributeTok{{-}ntsc} \AttributeTok{{-}nopatch} \AttributeTok{{-}verax16} \AttributeTok{{-}verax16{-}rom}\NormalTok{ vera\_pbi\_handler.rom }\DataTypeTok{\textbackslash{}}
     \AttributeTok{{-}run}\NormalTok{ TESTIRQ.COM}
\end{Highlighting}
\end{Shaded}

Verifica il gancio su \texttt{VIMIRQ}: VSYNC a 59,94 Hz, IRQ di riga,
mascheramento di AFLOW, catena verso l'OS e rimozione. Risultato atteso:
\textbf{PASS 13, FAIL 0}, sia in PAL sia in NTSC. Vedi
\texttt{Documentation/VERA-IRQ.md}.

\subsection{TESTPLR --- player VTM}\label{testplr-player-vtm}

\begin{Shaded}
\begin{Highlighting}[]
\VariableTok{$EMU} \AttributeTok{{-}xl} \AttributeTok{{-}pal} \AttributeTok{{-}nopatch} \AttributeTok{{-}volume}\NormalTok{ 100 }\AttributeTok{{-}verax16} \AttributeTok{{-}verax16{-}rom}\NormalTok{ vera\_pbi\_handler.rom }\DataTypeTok{\textbackslash{}}
\NormalTok{     disk3{-}standalone.atr}
\end{Highlighting}
\end{Shaded}

Dal DUP: \texttt{L}, \texttt{TESTPLR.COM}. Usare \texttt{-pal} o
\texttt{-ntsc} secondo la frequenza per cui è stato convertito il brano.
\texttt{-volume\ 100} serve: il PSG della VERA è basso al livello di
mixer predefinito (oppure alzare \texttt{-verax16-psg-volume}).

\subsection{TEST / TESTGS / TESTMAZ / TESTMTX --- driver e modi
video}\label{test-testgs-testmaz-testmtx-driver-e-modi-video}

\begin{Shaded}
\begin{Highlighting}[]
\VariableTok{$EMU} \AttributeTok{{-}xl} \AttributeTok{{-}pal} \AttributeTok{{-}nopatch} \AttributeTok{{-}verax16} \AttributeTok{{-}verax16{-}rom}\NormalTok{ vera\_pbi\_handler.rom }\DataTypeTok{\textbackslash{}}
\NormalTok{     disk2{-}veratests{-}80x30.atr}
\end{Highlighting}
\end{Shaded}

Usare il disco della risoluzione da provare (\texttt{40x30},
\texttt{80x30}, \texttt{80x60}) e avviare dal DUP \texttt{TEST8.COM},
\texttt{TESTGS8.COM}, \texttt{TESTMAZ8.COM}, \texttt{TESTMTX8.COM}
(cifra 4 o 6 sugli altri dischi): caricamento font, gradiente,
scrolling, demo labirinto e matrice.

\subsection{TESTRMT --- player RMT su POKEY +
VERA}\label{testrmt-player-rmt-su-pokey-vera}

Su tutti i dischi \texttt{disk2-veratests-*.atr}, non richiede
\texttt{VERA.SYS}. Dal DUP: \texttt{L}, \texttt{TESTRMT.COM}.

\begin{longtable}[]{@{}ll@{}}
\toprule\noalign{}
Tasto & Uscita \\
\midrule\noalign{}
\endhead
\bottomrule\noalign{}
\endlastfoot
\texttt{1} & solo POKEY \\
\texttt{2} & solo VERA \\
\texttt{3} & entrambi (default) \\
\texttt{4} & ibrido, solo RMT8: POKEY canali 1-4, VERA canali 5-8 \\
\texttt{S} & stereo VERA on/off \\
\texttt{ESC} & ferma ed esce \\
\end{longtable}

\subsection{TESTRIO --- musica durante l'I/O su
disco}\label{testrio-musica-durante-lio-su-disco}

\begin{Shaded}
\begin{Highlighting}[]
\FunctionTok{make}\NormalTok{ disk4{-}rmtio.atr}
\VariableTok{$EMU} \AttributeTok{{-}xl} \AttributeTok{{-}pal} \AttributeTok{{-}nopatchall} \AttributeTok{{-}nobasic} \DataTypeTok{\textbackslash{}}
     \AttributeTok{{-}verax16} \AttributeTok{{-}verax16{-}rom}\NormalTok{ vera\_pbi\_handler.rom disk4{-}rmtio.atr}
\end{Highlighting}
\end{Shaded}

Il disco avvia da solo \texttt{TESTRIO.COM} (autorun MyPicoDos).
\texttt{-nopatchall} è obbligatorio: senza, l'emulatore intercetta il
SIO e non c'è temporizzazione seriale reale. Atteso (PAL, loader RBL):
\textasciitilde720 byte/s, 0 errori, 0 tick persi.

\subsection{RUNCPM --- terminale ANSI tramite
FujiNet}\label{runcpm-terminale-ansi-tramite-fujinet}

Terminale 1, FujiNet-PC:

\begin{Shaded}
\begin{Highlighting}[]
\BuiltInTok{cd}\NormalTok{ FujiNet/fujinet{-}pc{-}ATARI}
\ExtensionTok{./run{-}fujinet} \AttributeTok{{-}c}\NormalTok{ fnconfig.ini }\AttributeTok{{-}s}\NormalTok{ SD/}
\end{Highlighting}
\end{Shaded}

Attendere \texttt{\#\#\#\ NetSIO\ stopped\ \#\#\#} nel log, poi nel
terminale 2:

\begin{Shaded}
\begin{Highlighting}[]
\VariableTok{$EMU} \AttributeTok{{-}xl} \AttributeTok{{-}pal} \AttributeTok{{-}nopatch} \AttributeTok{{-}netsio} \AttributeTok{{-}verax16} \AttributeTok{{-}verax16{-}rom}\NormalTok{ vera\_pbi\_handler.rom}
\end{Highlighting}
\end{Shaded}

Dettagli in \texttt{README-fujinet.md}.

\section{Risoluzione dei problemi}\label{risoluzione-dei-problemi}

\begin{longtable}[]{@{}
  >{\raggedright\arraybackslash}p{(\columnwidth - 2\tabcolsep) * \real{0.5000}}
  >{\raggedright\arraybackslash}p{(\columnwidth - 2\tabcolsep) * \real{0.5000}}@{}}
\toprule\noalign{}
\begin{minipage}[b]{\linewidth}\raggedright
Sintomo
\end{minipage} & \begin{minipage}[b]{\linewidth}\raggedright
Causa / rimedio
\end{minipage} \\
\midrule\noalign{}
\endhead
\bottomrule\noalign{}
\endlastfoot
\texttt{ERROR:\ VeraX16\ PBI\ card\ not\ found\ (\$D100\ not\ responding).}
& manca \texttt{-verax16}: ogni test chiama \texttt{vera\_require()}
all'avvio \\
Tre beep dall'altoparlante della console all'avvio & la ROM non ha visto
la VERA entro \textasciitilde0,7 s; controllare
\texttt{-verax16-config-ms} \\
Opzione \texttt{-verax16} sconosciuta & emulatore compilato senza
\texttt{-\/-enable-pbi-verax16} \\
VERA muta o molto bassa & aggiungere \texttt{-volume\ 100} o alzare
\texttt{-verax16-psg-volume} \\
Un file è sul disco ma il DOS non lo elenca, o ``0 FREE SECTORS'' & bug
di \texttt{dir2atr}; ricompilare con \texttt{make}, che esegue
\texttt{vera-tests/tools/fix\_atr\_vtoc.py} \\
Errori SIO o velocità sbagliata in \texttt{TESTRIO} & manca
\texttt{-nopatchall} \\
\end{longtable}

\end{document}
